1. Core Definitions
A Riesz space is a real vector space $E$ equipped with a partial order $\leq$ that is compatible with the vector structure and forms a lattice.

Compatibility Axioms
For all $x, y, z \in E$ and $\lambda \in \mathbb{R}_{\geq 0}$:
- Translation Invariance: If $x \leq y$, then $x + z \leq y + z$.
- Positive Homogeneity: If $x \leq y$, then $\lambda x \leq \lambda y$.
- Lattice Property: For every pair $x, y$, the supremum $x \vee y$ and infimum $x \wedge y$ exist in $E$.

The "Big Three" OperationsEvery element $x$ can be decomposed into its functional components:
- Positive part: $x^+ = x \vee 0$
- Negative part: $x^- = (-x) \vee 0$
- Absolute value: $|x| = x \vee (-x)$
- The Jordan Decomposition: $x = x^+ - x^-$ and $|x| = x^+ + x^-$.2. 

Key Substructures:
Understanding how to "slice" a Riesz space is vital for functional analysis.

Riesz Subspace subspace $L \subseteq E$ such that $x, y \in L \implies x \vee y \in L$.
Solid Set: Set $S$ where $Ideal$ solid subspace. (The kernels of lattice homomorphisms).

5. Orthogonality
Two elements $x, y$ are disjoint (denoted $x \perp y$) if $|x| \wedge |y| = 0$.
Disjoint Complement: $A^d = \{x \in E : |x| \wedge |a| = 0 \text{ for all } a \in A\}$. $A^d$ is always a band.

Band: An ideal $B$ that is "dedekind closed": if $D \subseteq B$ and $\sup D = x$, then $x \in B$.3. 
Important Types of Riesz Spaces: 
Notion of "completeness" the space is:
Archimedean: If $nx \leq y$ for all $n \in \mathbb{N}$ and $x \geq 0$, then $x = 0$. (No "infinitely small" elements).
Dedekind Complete: Every non-empty subset bounded above has a supremum.
Dedekind $\sigma$-complete: Every countable subset bounded above has a supremum.4. 

Fundamental Theorems

The Riesz Decomposition Property
If $0 \leq x \leq y_1 + y_2$, then there exist $x_1, x_2$ such that $x = x_1 + x_2$ with $0 \leq x_i \leq y_i$. 
Freudenthal’s Spectral Theorem:
In a Dedekind complete Riesz space with a principal unit $e$, every element $x$ in the principal ideal generated by $e$ can be approximated uniformly by "step functions" 
(linear combinations of projections of $e$ onto bands).

Kakutani’s Representation Theorem
Every AM-space (an $L^\infty$-like space) with a unit is lattice-isomorphic to $C(K)$, the space of continuous functions on a compact Hausdorff space $K$.

Riesz Projection Theorem: If $B$ is a projection band, then $E = B \oplus B^d$.6. 

Common Examples$C[0,1]$: Archimedean, but not Dedekind complete (pointwise suprema of continuous functions aren't always continuous).
$L^p(\mu)$: Dedekind complete for $1 \leq p \leq \infty$.
$\mathbb{R}^n$: With the component-wise ordering.
